\begin{pastexam}[label={px:ust-111-ee-ds-01}]{臺聯大電機類 111 學年度 第 1 題}
（設計）

\textbf{問答題}

(10\%)Please write a function to delete a selected node in a given singly linked list. For example,
given a linked list: 1->3->4->2 and the selected node of 4, your function should return the list as
1->3->2.
\end{pastexam}

\begin{pastsoln}
題目的「selected node of 4」可以讀成「值是 $4$ 的節點」或「指向那個節點的指標」。 下面先寫
給值的版本（最常見的讀法），再說明只給節點指標時怎麼做。 節點有 $\mathit{key}$、$\mathit{next}$
兩個欄位，函式回傳刪除後串列的第一個節點。

\textbf{給值 $k$，刪掉第一個值是 $k$ 的節點。} 單向串列要刪一個節點，必須改它\textbf{前一個}節點的
$\mathit{next}$（資結書第~5.1~節、虛擬碼~5.3），所以走的時候同時記著前一個節點 $p$。 要刪的是第一個
節點時沒有前一個，改 $\mathit{head}$。

\begin{pseudo}{刪除值為 $k$ 的節點}
\Function{Delete-Key}{$\mathit{head}, k$}
    \State $p \gets \NIL$；$x \gets \mathit{head}$
    \While{$x \ne \NIL$ \AND $x.\mathit{key} \ne k$}
        \State $p \gets x$；$x \gets x.\mathit{next}$
    \EndWhile
    \If{$x = \NIL$}
        \State \Return $\mathit{head}$ \Comment{沒有這個值，串列不變}
    \EndIf
    \If{$p = \NIL$}
        \State $\mathit{head} \gets x.\mathit{next}$ \Comment{刪的是第一個節點}
    \Else
        \State $p.\mathit{next} \gets x.\mathit{next}$
    \EndIf
    \State 釋放 $x$
    \State \Return $\mathit{head}$
\EndFunction
\end{pseudo}

\textbf{正確性。} 迴圈結束時，$x$ 是第一個值為 $k$ 的節點（或 \NIL），而且 $p$ 一直是 $x$ 的前一個
節點（$x$ 是第一個時 $p = \NIL$）。 讓前一個節點（或 $\mathit{head}$）直接指向 $x.\mathit{next}$，
$x$ 就不在串列裡了，其他節點的順序不變。

\textbf{例子}：$1 \to 3 \to 4 \to 2$，$k = 4$。 走到 $x = 4$ 時 $p = 3$，令 $3.\mathit{next} \gets 2$，
串列變成 $1 \to 3 \to 2$，回傳的 $\mathit{head}$ 仍是 $1$。

\textbf{複雜度}：最壞走過整條串列，$O(n)$；額外空間 $O(1)$。

\textbf{只給節點指標 $x$（拿不到前一個節點）時。} 單向串列從 $x$ 往回找不到前一個節點。 如果 $x$
不是最後一個節點，可以把下一個節點的內容抄到 $x$，再刪掉下一個節點：$y \gets x.\mathit{next}$，
$x.\mathit{key} \gets y.\mathit{key}$，$x.\mathit{next} \gets y.\mathit{next}$，釋放 $y$，$O(1)$。
例子裡 $x$ 是 $4$ 那個節點：抄進 $2$、拿掉原本的 $2$，得到 $1 \to 3 \to 2$。 這招對最後一個節點沒用
（沒有下一個可以抄），而且外面如果有指標指著原本的 $y$ 會失效，這時只能從頭走，回到上面的
做法。
\end{pastsoln}
